% ============================================================
% DROP-IN BLOCK: iVAE IDENTIFIABILITY THEORY
% Insert after \subsection{Structured disentangled generative models}
% (currently §2.3 in main.tex, before \subsection{Grokking as a phase transition})
% ============================================================

\subsection{Identifiability of the causal latent factor}
\label{sec:identifiability}

A central concern for any latent-variable method is \emph{identifiability}: 
does the encoder recover the \emph{true} latent factors, or some 
reparameterization that achieves good reconstruction without causal meaning?
We show that the Structured pi-SAE satisfies the conditions of iVAE
\citep{khemakhem2020variational}, which gives a formal guarantee that 
$z_\text{causal}$ recovers the true causal variable up to a 
component-wise transformation.

\begin{assumption}[Sufficient conditions for iVAE identifiability]
\label{ass:ivae}
\leavevmode
\begin{enumerate}[leftmargin=*,label=(\roman*)]
    \item \textbf{Auxiliary variable.} There exists an observed auxiliary
    variable $u$ (here, the task label $y$) such that the true latent prior
    factors as $p(z \mid u)$, with $p(z)$ not identifying $z$.
    \item \textbf{Label-conditional prior.} The conditional prior
    $p(z_\text{causal} \mid y)$ is an exponential family distribution whose
    sufficient statistics $T_j(z_j)$ are differentiable and whose natural
    parameters $\lambda_j(y)$ vary sufficiently across labels: for any
    $k$ latent dimensions, there exist $k+1$ labels $y_0, \ldots, y_k$ such
    that the matrix $(\lambda(y_1) - \lambda(y_0), \ldots,
    \lambda(y_k) - \lambda(y_0))$ has rank $k$.
    \item \textbf{Injective mixing function.} The decoder
    $g_\theta: \mathcal{Z} \to \mathcal{H}$ (mapping latent codes back to
    activation space $\mathcal{H} = \mathbb{R}^d$) is injective and
    differentiable.
\end{enumerate}
\end{assumption}

\begin{proposition}[Causal variable recovery]
\label{prop:recovery}
Let $h = g_\theta(z_\text{causal}, z_\text{nuisance})$ be activations
generated by an encoder satisfying Assumption~\ref{ass:ivae}.  Suppose the
Structured pi-SAE is trained to a global optimum of
$\mathcal{L}$ \emph{(Eq.~\ref{eq:vaeloss})} with $\beta > 0$ and $\alpha > 0$.
Then the learned encoder $\hat{q}(z_\text{causal} \mid h, y)$ recovers the
true causal variable up to a component-wise reparameterization:
\[
  \hat{z}_\text{causal} = \phi(z_\text{causal}),
\]
where $\phi: \mathbb{R}^k \to \mathbb{R}^k$ is a bijection acting
independently on each component (i.e., $\phi_j$ depends only on
$z_j$).
\end{proposition}

\begin{proof}[Proof sketch]
The ELBO term of $\mathcal{L}$ is tight at a global optimum, which forces
$q(z \mid h, y) = p(z \mid h, y)$ (the true posterior).  The KL term with
the structured prior $p(z_\text{causal} \mid y)$ forces the
encoder to match the label-conditional distribution.  Theorem~1 of
\citet{khemakhem2020variational} then guarantees that any two encoders
achieving the same ELBO value are related by a component-wise bijection on
$z_\text{causal}$, under Assumptions~(i)--(iii).  The supervised
classification term ($\alpha \mathcal{L}_\text{CE}$) additionally pins the
label-predictive direction of $z_\text{causal}$, removing the
permutation ambiguity. \qed
\end{proof}

\paragraph{Implications for causal abstraction.}
Proposition~\ref{prop:recovery} has three consequences.
\textbf{(1) NL-DAS cannot satisfy Assumption~(iii)} because its
unconstrained decoder is not forced to be injective; the training objective
provides no incentive for $g$ to be a bijection, so multiple latent codes
can map to the same activation (\S\ref{sec:nldas_vacuous}).
\textbf{(2) The component-wise ambiguity is harmless for IIA.}
Interchange interventions swap $z_\text{causal}$ wholesale; any
component-wise bijection $\phi$ is undone by the decoder, leaving IIA
invariant.  What matters for IIA is that $z_\text{causal}$ is
informationally equivalent to the true causal variable---and
Proposition~\ref{prop:recovery} guarantees exactly this.
\textbf{(3) Linear DAS corresponds to $\phi = $ rotation.}  When the true
causal variable is linear ($z_\text{causal} = Q^\top h$ for some orthogonal
$Q$), the structured encoder degenerates to a rotation, matching the DAS
solution.  Identifiability ``up to component-wise bijection'' reduces to
identifiability ``up to rotation in $O(k)$''---the standard DAS gauge
freedom.

\paragraph{Why the structured prior matters.}
Without the label-conditional prior (i.e., plain $\mathcal{N}(0,I)$),
Assumption~(i) fails: $p(z)$ does not depend on the auxiliary variable $y$,
and the iVAE theorem does not apply.  Empirically, plain-prior VAEs achieve
$\text{IIA} \approx 0$ on grokked operations despite high reconstruction
quality (Table~\ref{tab:2x2_ablation}), confirming that reconstruction alone
does not force causal recovery.  The structured prior is the mechanism by
which the label information ``pins'' the latent space to the correct
causal direction.

% NOTE: The full iVAE Theorem 1 statement and its exact conditions are in
% Khemakhem et al. 2020, Theorem 1 (p. 4). The exponential family
% formulation is Eq. (6) in that paper. Our prior p(z_causal | y) is an
% isotropic Gaussian N(mu_y, sigma_y * I) parameterized by the label;
% this satisfies the exponential family condition with sufficient statistics
% T_j(z_j) = (z_j, z_j^2) and natural parameters lambda_j(y) =
% (mu_y,j / sigma_y^2, -1/(2*sigma_y^2)).  The variability condition
% (rank condition on lambda differences) holds whenever the per-label means
% mu_y are sufficiently distinct -- satisfied in practice since each
% label corresponds to a distinct modular arithmetic output.

